% Options for packages loaded elsewhere
% Options for packages loaded elsewhere
\PassOptionsToPackage{unicode}{hyperref}
\PassOptionsToPackage{hyphens}{url}
\PassOptionsToPackage{dvipsnames,svgnames,x11names}{xcolor}
%
\documentclass[
  letterpaper,
]{article}
\usepackage{xcolor}
\usepackage{amsmath,amssymb}
\setcounter{secnumdepth}{5}
\usepackage{iftex}
\ifPDFTeX
  \usepackage[T1]{fontenc}
  \usepackage[utf8]{inputenc}
  \usepackage{textcomp} % provide euro and other symbols
\else % if luatex or xetex
  \usepackage{unicode-math} % this also loads fontspec
  \defaultfontfeatures{Scale=MatchLowercase}
  \defaultfontfeatures[\rmfamily]{Ligatures=TeX,Scale=1}
\fi
\usepackage{lmodern}
\ifPDFTeX\else
  % xetex/luatex font selection
\fi
% Use upquote if available, for straight quotes in verbatim environments
\IfFileExists{upquote.sty}{\usepackage{upquote}}{}
\IfFileExists{microtype.sty}{% use microtype if available
  \usepackage[]{microtype}
  \UseMicrotypeSet[protrusion]{basicmath} % disable protrusion for tt fonts
}{}
\makeatletter
\@ifundefined{KOMAClassName}{% if non-KOMA class
  \IfFileExists{parskip.sty}{%
    \usepackage{parskip}
  }{% else
    \setlength{\parindent}{0pt}
    \setlength{\parskip}{6pt plus 2pt minus 1pt}}
}{% if KOMA class
  \KOMAoptions{parskip=half}}
\makeatother
% Make \paragraph and \subparagraph free-standing
\makeatletter
\ifx\paragraph\undefined\else
  \let\oldparagraph\paragraph
  \renewcommand{\paragraph}{
    \@ifstar
      \xxxParagraphStar
      \xxxParagraphNoStar
  }
  \newcommand{\xxxParagraphStar}[1]{\oldparagraph*{#1}\mbox{}}
  \newcommand{\xxxParagraphNoStar}[1]{\oldparagraph{#1}\mbox{}}
\fi
\ifx\subparagraph\undefined\else
  \let\oldsubparagraph\subparagraph
  \renewcommand{\subparagraph}{
    \@ifstar
      \xxxSubParagraphStar
      \xxxSubParagraphNoStar
  }
  \newcommand{\xxxSubParagraphStar}[1]{\oldsubparagraph*{#1}\mbox{}}
  \newcommand{\xxxSubParagraphNoStar}[1]{\oldsubparagraph{#1}\mbox{}}
\fi
\makeatother


\usepackage{longtable,booktabs,array}
\usepackage{calc} % for calculating minipage widths
% Correct order of tables after \paragraph or \subparagraph
\usepackage{etoolbox}
\makeatletter
\patchcmd\longtable{\par}{\if@noskipsec\mbox{}\fi\par}{}{}
\makeatother
% Allow footnotes in longtable head/foot
\IfFileExists{footnotehyper.sty}{\usepackage{footnotehyper}}{\usepackage{footnote}}
\makesavenoteenv{longtable}
\usepackage{graphicx}
\makeatletter
\newsavebox\pandoc@box
\newcommand*\pandocbounded[1]{% scales image to fit in text height/width
  \sbox\pandoc@box{#1}%
  \Gscale@div\@tempa{\textheight}{\dimexpr\ht\pandoc@box+\dp\pandoc@box\relax}%
  \Gscale@div\@tempb{\linewidth}{\wd\pandoc@box}%
  \ifdim\@tempb\p@<\@tempa\p@\let\@tempa\@tempb\fi% select the smaller of both
  \ifdim\@tempa\p@<\p@\scalebox{\@tempa}{\usebox\pandoc@box}%
  \else\usebox{\pandoc@box}%
  \fi%
}
% Set default figure placement to htbp
\def\fps@figure{htbp}
\makeatother





\setlength{\emergencystretch}{3em} % prevent overfull lines

\providecommand{\tightlist}{%
  \setlength{\itemsep}{0pt}\setlength{\parskip}{0pt}}



 


\usepackage{orcidlink}  
\usepackage[nonatbib, final]{neurips_2023}
\usepackage{amsmath}
\usepackage{amsthm}
\usepackage[utf8]{inputenc} % allow utf-8 input
\usepackage[T1]{fontenc}    % use 8-bit T1 fonts
\usepackage{hyperref}       % hyperlinks
\usepackage{url}            % simple URL typesetting
\usepackage{booktabs}       % professional-quality tables
\usepackage{amsfonts}       % blackboard math symbols
\usepackage{nicefrac}       % compact symbols for 1/2, etc.
\usepackage{microtype}      % microtypography
\usepackage{xcolor}         % colors
\usepackage[numbers]{natbib}
\usepackage{comment}
\usepackage{graphicx} 
\bibliographystyle{abbrvnat}
\makeatletter
\@ifpackageloaded{caption}{}{\usepackage{caption}}
\AtBeginDocument{%
\ifdefined\contentsname
  \renewcommand*\contentsname{Table of contents}
\else
  \newcommand\contentsname{Table of contents}
\fi
\ifdefined\listfigurename
  \renewcommand*\listfigurename{List of Figures}
\else
  \newcommand\listfigurename{List of Figures}
\fi
\ifdefined\listtablename
  \renewcommand*\listtablename{List of Tables}
\else
  \newcommand\listtablename{List of Tables}
\fi
\ifdefined\figurename
  \renewcommand*\figurename{Figure}
\else
  \newcommand\figurename{Figure}
\fi
\ifdefined\tablename
  \renewcommand*\tablename{Table}
\else
  \newcommand\tablename{Table}
\fi
}
\@ifpackageloaded{float}{}{\usepackage{float}}
\floatstyle{ruled}
\@ifundefined{c@chapter}{\newfloat{codelisting}{h}{lop}}{\newfloat{codelisting}{h}{lop}[chapter]}
\floatname{codelisting}{Listing}
\newcommand*\listoflistings{\listof{codelisting}{List of Listings}}
\makeatother
\makeatletter
\makeatother
\makeatletter
\@ifpackageloaded{caption}{}{\usepackage{caption}}
\@ifpackageloaded{subcaption}{}{\usepackage{subcaption}}
\makeatother
\usepackage{bookmark}
\IfFileExists{xurl.sty}{\usepackage{xurl}}{} % add URL line breaks if available
\urlstyle{same}
\hypersetup{
  pdftitle={Comparing the Impact and Arrival Cost of Naive and OBI-Enhanced Order Splitting Strategies},
  pdfauthor={Michael Cusack-Nelkin; Alick Zhou; Debby Zhao; Michael Bakshandeh; Vijeet Prasad},
  pdfkeywords={optimal, execution, order, splitting, imbalance},
  colorlinks=true,
  linkcolor={blue},
  filecolor={Maroon},
  citecolor={Blue},
  urlcolor={Blue},
  pdfcreator={LaTeX via pandoc}}



\title{Comparing the Impact and Arrival Cost of Naive and OBI-Enhanced
Order Splitting Strategies}
\author{
\textbf{Michael Cusack-Nelkin}%
\\%
%
\texttt{mc5882@columbia.edu}\and 
\textbf{Alick Zhou}%
\\%
%
\texttt{az2960@columbia.edu}\and 
\textbf{Debby Zhao}%
\\%
%
\texttt{lz3083@columbia.edu}\and 
\textbf{Michael Bakshandeh}%
\\%
%
\texttt{mb5499@columbia.edu}\and 
\textbf{Vijeet Prasad}%
\\%
%
\texttt{vp2580@columbia.edu}}
\date{2026-04-23}
\begin{document}
\maketitle


\section{Motivation and Main Ideas}\label{motivation-and-main-ideas}

Motivation is a way to minimize the cost of an order by taking advantage
of favorable price moves and avoiding unfavorable ones.

\subsection{Definitions}\label{definitions}

\begin{itemize}
\tightlist
\item
  OBi
\item
  Impact
\item
  Arrival Cost
\end{itemize}

\subsection{Naive Splitting is Bad}\label{naive-splitting-is-bad}

Moves prices too much
\pandocbounded{\includegraphics[keepaspectratio]{.pdf}}

\subsection{OBI Provides a Way to Identify Favorable/Adverse Price
Movements}\label{obi-provides-a-way-to-identify-favorableadverse-price-movements}

OBI of the top 6 levels.
\pandocbounded{\includegraphics[keepaspectratio]{.pdf}}

\section{Empirical Findings}\label{empirical-findings}

\subsection{Descriptive Stats}\label{descriptive-stats}

\subsection{Test Stats / Significance}\label{test-stats-significance}

\subsection{Latency / Replenishment / Time of Day
Plots/Stats}\label{latency-replenishment-time-of-day-plotsstats}

\pandocbounded{\includegraphics[keepaspectratio]{.pdf}}
\pandocbounded{\includegraphics[keepaspectratio]{.pdf}}
\pandocbounded{\includegraphics[keepaspectratio]{.pdf}}

\section{Methodological Choices and
Innovations}\label{methodological-choices-and-innovations}

\subsection{Failed Attempt: OBI Aggressive
Only}\label{failed-attempt-obi-aggressive-only}

Shifting reference frame from alpha to execution (given mandate to
buy/sell.) Not trying to predict price, but take advantage of price
moves. \pandocbounded{\includegraphics[keepaspectratio]{.pdf}}

\subsection{Successful: OBI Calibrated
Aggression}\label{successful-obi-calibrated-aggression}

\begin{figure}

\centering{

\pandocbounded{\includegraphics[keepaspectratio]{.pdf}}

}

\caption{\label{fig-hybrid}Hybrid Plot}

\end{figure}%

\subsection{Simulator}\label{simulator}

\begin{figure}

\centering{

\pandocbounded{\includegraphics[keepaspectratio]{.pdf}}

}

\caption{\label{fig-sim}Impact Plot}

\end{figure}%

\subsubsection{Choice: fixed horizon (time
constraint)}\label{choice-fixed-horizon-time-constraint}

\subsubsection{Choice: constant fills
within-band}\label{choice-constant-fills-within-band}

\section{Future Work}\label{future-work}

\subsection{Smoothed OBI}\label{smoothed-obi}

Tried it, got some weird results (no hyperparam tuning, possibly a quirk
of replenishment.)

\subsection{A Catch-up/Fall-Behind
Mechanism}\label{a-catch-upfall-behind-mechanism}

Currently filling naively within-band.




\end{document}
